% Source: Gerzensee "Real Analysis" slides 58-76 (Cauchy sequences, completeness,
% the (0,1] counterexample, products and closed subsets, the contraction mapping
% theorem with its four-step proof and error bound, the list of applications, the
% sup norm on B(R^n,R), Blackwell's sufficient conditions and their proof).
% Supplementary: Fukuda (2026), Sec. 8.4 (pp. 278-294), including Prop. 8.20
% (subspaces of R^n), Prop. 8.21 (Blackwell), Exercises 8.26-8.31 and Remark 8.11.
% External references actually cited in this chapter: Lucas (1978), Stokey et
% al. (1989).

\chapter{Completeness, Contractions, and Fixed Points}
\label{ch:completeness}

\section{Motivation: Proving a Limit Exists Without Knowing It}\label{sec:comp-motivation}

The definition of convergence names the limit. To use it one must already have a
candidate. In most of the arguments that matter --- iterating a Bellman
operator, solving a differential equation, constructing an equilibrium price ---
the object whose existence is at issue is precisely the limit, and there is no
candidate to name.

Cauchy's condition removes the limit from the statement: it asks only that the
terms of the sequence eventually be close \emph{to one another}. In a complete
space that internal condition is enough. Completeness is thus the hypothesis
that converts a self-referential existence claim into a verifiable one, and it
is the hypothesis on which the contraction mapping theorem, and with it the
entire recursive method in macroeconomics, rests.

\section{Cauchy Sequences}\label{sec:comp-cauchy}

\begin{definition}[Cauchy sequence]\label{def:cauchy}
	A sequence $(x_n)_{n\in\N}$ in a metric space $(X,d)$ is a
	\dfn{Cauchy sequence} if for every $\varepsilon>0$ there is $n_0\in\N$ such that
	$m,n\geq n_0$ implies $d(x_m,x_n)<\varepsilon$.
\end{definition}

\begin{proposition}[Convergent implies Cauchy]\label{prop:convergent-cauchy}
	Every convergent sequence in a metric space is Cauchy.
\end{proposition}

\begin{proof}
	Let $x_n\to x^\ast$ and $\varepsilon>0$. Choose $n_0$ with
	$d(x^\ast,x_n)<\varepsilon/2$ for $n\geq n_0$. For $m,n\geq n_0$,
	$d(x_m,x_n)\leq d(x_m,x^\ast)+d(x^\ast,x_n)<\varepsilon$.
\end{proof}

\begin{proposition}[Cauchy implies bounded]\label{prop:cauchy-bounded}
	Every Cauchy sequence in a metric space is bounded.
\end{proposition}

\begin{proof}
	Take $\varepsilon=1$ and the corresponding $n_0$, so $d(x_{n_0},x_n)<1$ for
	$n\geq n_0$. Put
	\[
		M\eqdef\max\bigl\{1,\ d(x_{n_0},x_1),\ \dots,\ d(x_{n_0},x_{n_0-1})\bigr\},
	\]
	a maximum over finitely many numbers. Then $(x_n)\subseteq\ball{M+1}{x_{n_0}}$.
\end{proof}

These are Propositions~8.14 and~8.15 of \textcite{fukuda2026notes}
\autocite[slides 58--59]{fukuda2026analysis}.

\begin{lemma}[Cauchy with a convergent subsequence]\label{lem:cauchy-subseq}
	If a Cauchy sequence has a subsequence converging to $x^\ast$, then the whole
	sequence converges to $x^\ast$.
\end{lemma}

\begin{proof}
	Let $\varepsilon>0$. Choose $n_0$ with $d(x_m,x_n)<\varepsilon/2$ for
	$m,n\geq n_0$, and choose an index $n_k\geq n_0$ of the subsequence with
	$d(x^\ast,x_{n_k})<\varepsilon/2$. For $n\geq n_0$,
	$d(x^\ast,x_n)\leq d(x^\ast,x_{n_k})+d(x_{n_k},x_n)<\varepsilon$.
\end{proof}

The converse of \autoref{prop:convergent-cauchy} is false in general: the
implication ``Cauchy $\Rightarrow$ convergent'' is not a theorem about metric
spaces but a property that some metric spaces have and others lack. That
property is completeness.

\section{Complete Metric Spaces}\label{sec:comp-complete}

\begin{definition}[Completeness]\label{def:complete-space}
	A metric space $(X,d)$ is \dfnas{complete metric space}{complete} if every
	Cauchy sequence in $X$ converges to a limit belonging to $X$. A subset
	$A\subseteq X$ is complete if $(A,\restr{d}{A\times A})$ is complete.
\end{definition}

In a complete space the two conditions coincide, and the practical consequence is
the one stated on \autocite[slide 60]{fukuda2026analysis}: to prove that
$(x_n)$ has a limit it suffices to prove that $(x_n)$ is Cauchy, a condition
involving only the terms of the sequence.

\begin{eg}[An incomplete space]\label{eg:incomplete-halfopen}
	The space $X=(0,1]$ with the metric inherited from $\R$ is not complete. The
	sequence $x_n=1/n$ is Cauchy --- given $\varepsilon>0$,
	\autoref{prop:archimedean} supplies $n_0$ with $1/n_0<\varepsilon$, and then
	$\abs{1/m-1/n}\leq1/n_0<\varepsilon$ for $m,n\geq n_0$ --- but its limit $0$ does
	not belong to $X$ \autocite[slide 61]{fukuda2026analysis}.

	Nothing about the sequence has changed relative to $\R$; what changed is the
	ambient set. Completeness is not a property of a sequence but of a space, and it
	is destroyed by deleting a single point.
\end{eg}

\begin{theorem}[$\R^n$ is complete]\label{thm:rn-complete}
	$\R$ is a complete metric space, and so is $\R^n$ for every finite $n$.
\end{theorem}

\begin{proof}
	Let $(x_k)$ be Cauchy in $\R$. By \autoref{prop:cauchy-bounded} it is bounded,
	so by \autoref{thm:bolzano-weierstrass} it has a convergent subsequence, and by
	\autoref{lem:cauchy-subseq} the whole sequence converges. The limit is a real
	number, so $\R$ is complete. The case of $\R^n$ is identical, using the
	$\R^n$ form of \autoref{thm:bolzano-weierstrass}.
\end{proof}

The chain of implications locates completeness within the logical structure of
the real numbers: the least upper
bound property (\autoref{ass:completeness-axiom}) gives monotone convergence
(\autoref{prop:monotone-convergence}), which with the monotone subsequence lemma
gives Bolzano--Weierstrass, which with \autoref{lem:cauchy-subseq} gives metric
completeness. Completeness of $\R$ as a metric space is therefore a theorem, not
an axiom, once the order-completeness axiom is in place.

\begin{proposition}[Products]\label{prop:product-complete}
	If $X$ and $Y$ are complete metric spaces, then $X\times Y$ with the metric
	$d\bigl((x,y),(x',y')\bigr)=\bigl(d_X(x,x')^2+d_Y(y,y')^2\bigr)^{1/2}$ is
	complete.
\end{proposition}

\begin{proof}
	If $(x_k,y_k)$ is Cauchy in the product, then $d_X(x_j,x_k)\leq d((x_j,y_j),(x_k,y_k))$
	shows $(x_k)$ is Cauchy in $X$, and symmetrically for $(y_k)$. Their limits
	$x^\ast,y^\ast$ exist by completeness, and
	$d((x_k,y_k),(x^\ast,y^\ast))\to0$.
\end{proof}

\begin{proposition}[Complete subsets are exactly the closed ones]
	\label{prop:closed-subset-complete}
	Let $X$ be a complete metric space and $A\subseteq X$. Then $A$ is complete if
	and only if $A$ is closed in $X$.
\end{proposition}

\begin{proof}
	($\Leftarrow$) Let $A$ be closed and $(x_n)\subseteq A$ Cauchy. By completeness
	of $X$ it converges to some $x^\ast\in X$, and by
	\autoref{thm:closed-sequential} $x^\ast\in A$. So the sequence converges in $A$.

	($\Rightarrow$) Let $A$ be complete and $x^\ast\in\closure A$. By the
	construction in the proof of \autoref{thm:closed-sequential} there is
	$(x_n)\subseteq A$ with $x_n\to x^\ast$; this sequence is Cauchy by
	\autoref{prop:convergent-cauchy}, hence converges to a point of $A$ by
	completeness of $A$, and by uniqueness of limits
	(\autoref{prop:sequence-basics}(i)) that point is $x^\ast$. So
	$\closure A\subseteq A$.
\end{proof}

This is Proposition~8.19 of \textcite{fukuda2026notes}
\autocite[slide 62]{fukuda2026analysis}. Note the asymmetry of the hypotheses:
completeness of the ambient space is used only in one direction, and the
equivalence is a statement about subsets of a complete space, not about metric
spaces in general.

\begin{corollary}\label{cor:subspace-complete}
	Every vector subspace of $\R^n$ is complete, and hence closed.
\end{corollary}

\begin{proof}
	Let $V\subseteq\R^n$ be a subspace with basis $v_1,\dots,v_k$ and let $(x_j)$ be
	Cauchy in $V$. Completing to a basis of $\R^n$ and using
	\autoref{thm:rn-complete}, $x_j\to x^\ast\in\R^n$. Writing
	$x_j=\sum_{i=1}^k c^i_jv_i$, the coordinate functionals are linear, hence
	continuous by \autoref{prop:linear-continuous}(ii), so the coefficient vectors
	$(c_j)$ converge in $\R^k$ to some $c^\ast$ and $x^\ast=\sum_ic^{i\ast}v_i\in V$.
	Closedness follows from \autoref{prop:closed-subset-complete}.
\end{proof}

This is Proposition~8.20 of \textcite[p.~279]{fukuda2026notes}. In infinite
dimensions the conclusion fails --- the subspace of finitely supported sequences
is not closed in $\ell^2$ --- and this is one of the practical differences that
\autoref{ch:projection} must accommodate: the projection theorem is stated for
closed subspaces because in infinite dimensions ``subspace'' does not imply
``closed''.

\subsection{Function spaces}

The space in which dynamic programming operates is not $\R^n$. It is a space of
functions, and its completeness is what the whole method depends on.

\begin{theorem}[The bounded functions are a Banach space]\label{thm:B-complete}
	Let $S$ be any non-empty set and let $\Bspace{S}{\R}$ be the set of bounded
	functions $f\colon S\to\R$, equipped with the supremum norm
	$\norm{f}_\infty=\sup_{s\in S}\abs{f(s)}$. Then $\Bspace{S}{\R}$ is a Banach
	space.
\end{theorem}

\begin{proof}
	That $\norm{\cdot}_\infty$ is a norm is \autoref{prop:sup-norm}. For
	completeness, let $(f_k)$ be Cauchy in $\Bspace{S}{\R}$.

	\emph{Step 1: pointwise limit.} For each fixed $s\in S$,
	$\abs{f_j(s)-f_k(s)}\leq\norm{f_j-f_k}_\infty$, so $(f_k(s))_k$ is Cauchy in
	$\R$ and converges by \autoref{thm:rn-complete}. Define
	$f(s)\eqdef\lim_kf_k(s)$.

	\emph{Step 2: the convergence is uniform.} Given $\varepsilon>0$, take $k_0$
	with $\norm{f_j-f_k}_\infty<\varepsilon/2$ for $j,k\geq k_0$. Fix $s$ and
	$k\geq k_0$, and let $j\to\infty$ in $\abs{f_j(s)-f_k(s)}<\varepsilon/2$; the
	absolute value is continuous, so $\abs{f(s)-f_k(s)}\leq\varepsilon/2$. The bound
	is uniform in $s$, so $\norm{f-f_k}_\infty\leq\varepsilon/2<\varepsilon$ for
	every $k\geq k_0$.

	\emph{Step 3: the limit is bounded.} With $\varepsilon=1$ and the corresponding
	$k_0$, $\norm{f}_\infty\leq\norm{f-f_{k_0}}_\infty+\norm{f_{k_0}}_\infty
		\leq1+\norm{f_{k_0}}_\infty<\infty$. Hence $f\in\Bspace{S}{\R}$ and
	$f_k\to f$ in the supremum norm.
\end{proof}

\begin{remark}[Supremum-norm convergence is uniform convergence]
	\label{rem:uniform-convergence}
	Step 2 is the content of the identification: $\norm{f-f_k}_\infty\to0$ says
	exactly that $f_k\to f$ uniformly on $S$, that is, with a single $k_0$ serving
	every $s$ simultaneously. Pointwise convergence alone is a strictly weaker
	statement and does not preserve continuity: $f_k(s)=s^k$ on $[0,1]$ converges
	pointwise to the discontinuous limit that is $0$ on $[0,1)$ and $1$ at $1$, and
	$\norm{f-f_k}_\infty=1$ for every $k$. The supremum norm is chosen for dynamic
	programming because convergence in it is strong enough to carry
	continuity, monotonicity and concavity to the limit, as
	\autoref{cor:invariant-set} exploits.
\end{remark}

\begin{proposition}[Bounded continuous functions]\label{prop:C-complete}
	Let $S\subseteq\R^n$. The set $\Cspace{S}{\R}\cap\Bspace{S}{\R}$ of bounded
	continuous functions is a closed subset of $\Bspace{S}{\R}$, hence a Banach
	space. If $S$ is compact, every continuous function on $S$ is bounded by
	\autoref{cor:weierstrass-continuous}, so $\Cspace{S}{\R}$ itself is a Banach
	space in the supremum norm.
\end{proposition}

\begin{proof}
	By \autoref{prop:closed-subset-complete} it suffices to show closedness. Let
	$f_k\to f$ uniformly with each $f_k$ continuous, fix $s_0\in S$ and
	$\varepsilon>0$. Choose $k$ with $\norm{f-f_k}_\infty<\varepsilon/3$ and then
	$\delta>0$ with $\abs{f_k(s)-f_k(s_0)}<\varepsilon/3$ for
	$s\in\ball{\delta}{s_0}\cap S$. For such $s$,
	\[
		\abs{f(s)-f(s_0)}\leq\abs{f(s)-f_k(s)}+\abs{f_k(s)-f_k(s_0)}
		+\abs{f_k(s_0)-f(s_0)}<\varepsilon .
	\]
\end{proof}

\subsection{Compactness in a function space}\label{subsec:comp-ascoli}

Completeness of $\Cspace{K}{\R}$ delivers limits of Cauchy sequences. It does not
deliver convergent subsequences of bounded sequences, and in an
infinite-dimensional space the Heine--Borel theorem is false: the closed unit
ball of $\Cspace{[0,1]}{\R}$ is closed and bounded and not compact. Take
$f_k(x)=x^{k}$, so that $\norm{f_k}_\infty=1$ for every $k$. Every subsequence
still converges pointwise to the function equal to $0$ on $[0,1)$ and to $1$ at
$x=1$, which is discontinuous; a uniform limit of continuous functions is
continuous by \autoref{prop:C-complete}, so no subsequence converges
uniformly. This is the same failure that
\autoref{eg:weierstrass-failures} exhibits for the Weierstrass theorem, and the
repair is the same: replace boundedness by a condition that controls the whole
family at once.

\begin{definition}[Equicontinuity]\label{def:equicontinuous}
	Let $(X,d)$ be a metric space. A family $\mathcal{F}\subseteq\Cspace{X}{\R}$ is
	\dfn{equicontinuous} if for every $\varepsilon>0$ there is $\delta>0$ such that
	\[
		d(x,y)<\delta\quad\Longrightarrow\quad\abs{f(x)-f(y)}<\varepsilon
		\qquad\text{for every }f\in\mathcal{F}.
	\]
	It is \dfnas{uniformly bounded family}{uniformly bounded} if
	$\sup_{f\in\mathcal{F}}\norm{f}_\infty<\infty$.
\end{definition}

The quantifier order carries the content: $\delta$ depends on $\varepsilon$ but
not on $f$, whereas continuity of each member separately would allow $\delta$ to
shrink as $f$ ranges over the family. The family $\{x\mapsto x^{k}\}_{k\in\N}$
fails the condition: for any $\delta\in(0,1)$ the pair $x=1$, $y=1-\delta$ gives
$\abs{f_k(x)-f_k(y)}=1-(1-\delta)^{k}\to1$, so no $\delta$ works for
$\varepsilon=\tfrac12$.

\begin{theorem}[Arzelà--Ascoli]\label{thm:arzela-ascoli}
	Let $K\subseteq\R^n$ be compact and $\mathcal{F}\subseteq\Cspace{K}{\R}$. Then
	$\mathcal{F}$ has compact closure in $(\Cspace{K}{\R},\norm{\cdot}_\infty)$ if
	and only if $\mathcal{F}$ is uniformly bounded and equicontinuous.
\end{theorem}

\begin{proof}
	\emph{Sufficiency.} Let $(f_k)$ be a sequence in $\mathcal{F}$. A compact subset
	of $\R^n$ has a countable dense subset $\{x^1,x^2,\dots\}$: for each $m$ cover
	$K$ by finitely many balls of radius $1/m$ and collect the centres. The
	sequence $(f_k(x^1))_k$ is bounded in $\R$, so by
	\autoref{thm:bolzano-weierstrass} it has a convergent subsequence
	$(f_{k}^{1})$; from it extract $(f_{k}^{2})$ converging at $x^2$, and so on.
	The diagonal sequence $g_j\eqdef f_{j}^{j}$ is, from index $m$ onwards, a
	subsequence of $(f_k^m)$, so $(g_j(x^m))_j$ converges for every $m$.

	Fix $\varepsilon>0$ and take $\delta$ from equicontinuity for $\varepsilon/3$.
	Compactness of $K$ gives finitely many points $x^{i_1},\dots,x^{i_p}$ of the
	dense set with $K\subseteq\bigcup_{l}\ball{\delta}{x^{i_l}}$. Choose $J$ so that
	$\abs{g_j(x^{i_l})-g_{j'}(x^{i_l})}<\varepsilon/3$ for all $j,j'\geq J$ and all
	$l\leq p$. Given $x\in K$, pick $l$ with $d(x,x^{i_l})<\delta$; then for
	$j,j'\geq J$,
	\[
		\abs{g_j(x)-g_{j'}(x)}
		\leq\abs{g_j(x)-g_j(x^{i_l})}+\abs{g_j(x^{i_l})-g_{j'}(x^{i_l})}
		+\abs{g_{j'}(x^{i_l})-g_{j'}(x)}<\varepsilon,
	\]
	the outer terms by equicontinuity. Hence $\norm{g_j-g_{j'}}_\infty\leq
		\varepsilon$ for $j,j'\geq J$, so $(g_j)$ is Cauchy and converges in
	$\Cspace{K}{\R}$ by \autoref{prop:C-complete}. Every sequence in $\mathcal{F}$
	therefore has a uniformly convergent subsequence. Passing to the closure
	costs one approximation: given $(h_j)$ in $\closure{\mathcal{F}}$, choose
	$f_j\in\mathcal{F}$ with $\norm{h_j-f_j}_\infty<1/j$, extract
	$f_{j_i}\to f$ in $\Cspace{K}{\R}$, and note that
	$\norm{h_{j_i}-f}_\infty\leq 1/j_i+\norm{f_{j_i}-f}_\infty\to0$ with
	$f\in\closure{\mathcal{F}}$. So $\closure{\mathcal{F}}$ is sequentially
	compact, hence compact by \autoref{thm:compact-seq-compact}.

	\emph{Necessity.} A compact set in a metric space is bounded, which gives
	uniform boundedness, and totally bounded, so for $\varepsilon>0$ there are
	$h_1,\dots,h_N\in\Cspace{K}{\R}$ with every $f\in\mathcal{F}$ within
	$\varepsilon/3$ of some $h_i$ in supremum norm. Each $h_i$ is uniformly
	continuous on the compact $K$ by \autoref{thm:heine-cantor}; let $\delta$
	be the smallest of the finitely many moduli for $\varepsilon/3$. For
	$d(x,y)<\delta$ and $f$ within $\varepsilon/3$ of $h_i$,
	\[
		\abs{f(x)-f(y)}\leq\abs{f(x)-h_i(x)}+\abs{h_i(x)-h_i(y)}
		+\abs{h_i(y)-f(y)}<\varepsilon. \qedhere
	\]
\end{proof}

\begin{corollary}[Uniformly Lipschitz families]\label{cor:lipschitz-compact}
	Let $K\subseteq\R^n$ be compact and let $\mathcal{F}\subseteq\Cspace{K}{\R}$
	satisfy $\abs{f(x)-f(y)}\leq L\norm{x-y}$ for every $f\in\mathcal{F}$ and all
	$x,y\in K$, with a single constant $L$, and $\abs{f(x_0)}\leq M$ for some fixed
	$x_0\in K$ and every $f$. Then $\mathcal{F}$ has compact closure in
	$\Cspace{K}{\R}$.
\end{corollary}

\begin{proof}
	Equicontinuity holds with $\delta=\varepsilon/L$. Uniform boundedness follows
	from $\abs{f(x)}\leq\abs{f(x_0)}+L\norm{x-x_0}\leq M+L\diam K$. Apply
	\autoref{thm:arzela-ascoli}.
\end{proof}

\begin{remark}[Where this is used]\label{rem:ascoli-uses}
	\autoref{cor:lipschitz-compact} is the compactness input to existence proofs in
	continuous time. In \autoref{ch:continuoustime} the admissible state
	trajectories of $\dot x=f(x,u)$ with controls taking values in a compact set
	and $f$ bounded on the relevant region are all Lipschitz on $[0,T]$ with the
	same constant $\sup\norm{f}$ and share the initial value $x(0)=x_0$. The family
	of admissible trajectories is therefore relatively compact in
	$\Cspace{[0,T]}{\R^{n}}$, and a maximising sequence has a uniformly convergent
	subsequence; \autoref{rem:ct-existence} states the conclusion this yields and
	the extra convexity hypothesis needed to keep the limit admissible.

	The theorem is also what makes a family of value functions or policy functions
	generated by varying a parameter relatively compact, which is the standard route
	to continuity of an equilibrium in a parameter when no explicit formula is
	available. Boundedness alone never suffices, as $\sin(k\cdot)$ shows.
\end{remark}

\section{The Contraction Mapping Theorem}\label{sec:comp-contraction}

\begin{definition}[Contraction]\label{def:contraction}
	Let $(X,d)$ be a metric space. A map $T\colon X\to X$ is a
	\dfnas{contraction mapping}{contraction} with \dfn{modulus} $\beta$ if there is
	$\beta\in[0,1)$ with
	\[
		d(Tx,Ty)\leq\beta\,d(x,y)\qquad\text{for all }x,y\in X .
	\]
	A point $x^\ast\in X$ with $Tx^\ast=x^\ast$ is a \dfnas{fixed
		point}{fixed point} of $T$.
	\nidx{T}{$T$, operator; $Tx$, its value at $x$}
\end{definition}

Two features of the definition deserve emphasis. The modulus must be strictly
less than $1$ and must be uniform over all pairs $x,y$. The weaker condition
$d(Tx,Ty)<d(x,y)$ for $x\neq y$ --- a \emph{contractive} map --- does not
suffice: on $X=[1,\infty)$ the map $Tx=x+1/x$ satisfies it, is a self-map of a
complete space, and has no fixed point, because the implicit modulus approaches
$1$ as $x\to\infty$.

\begin{theorem}[Banach contraction mapping theorem]\label{thm:banach-fixed-point}
	Let $(X,d)$ be a non-empty complete metric space and let $T\colon X\to X$ be a
	contraction with modulus $\beta\in[0,1)$. Then $T$ has exactly one fixed point
	$x^\ast\in X$. Moreover, for any $x_0\in X$ the sequence of iterates
	$x_n\eqdef Tx_{n-1}=T^nx_0$ converges to $x^\ast$, and
	\begin{equation}\label{eq:contraction-error}
		d(x_n,x^\ast)\leq\frac{\beta^n}{1-\beta}\,d(x_0,Tx_0)
		\qquad\text{for every }n\in\N_0 .
	\end{equation}
\end{theorem}

\begin{proof}
	\emph{Step 1: the iterates are Cauchy.} Fix $x_0\in X$ and define
	$x_n=Tx_{n-1}$. Applying the contraction property $k$ times,
	$d(x_k,x_{k+1})=d(Tx_{k-1},Tx_k)\leq\beta\,d(x_{k-1},x_k)\leq\dots\leq
		\beta^k d(x_0,x_1)$. For $m>n$, the triangle inequality applied $m-n-1$ times
	gives
	\[
		d(x_n,x_m)\leq\sum_{k=n}^{m-1}d(x_k,x_{k+1})
		\leq\sum_{k=n}^{m-1}\beta^kd(x_0,x_1)
		=\beta^n\,\frac{1-\beta^{m-n}}{1-\beta}\,d(x_0,x_1)
		\leq\frac{\beta^n}{1-\beta}\,d(x_0,x_1),
	\]
	where the geometric sum is finite because $\beta<1$. Given $\varepsilon>0$,
	since $\beta^n\to0$ there is $n_0$ with
	$\beta^{n}d(x_0,x_1)/(1-\beta)<\varepsilon$ for $n\geq n_0$, and then
	$d(x_n,x_m)<\varepsilon$ for all $m,n\geq n_0$. So $(x_n)$ is Cauchy.

	\emph{Step 2: the limit exists.} $X$ is complete, so
	$x^\ast\eqdef\lim_nx_n$ exists in $X$.

	\emph{Step 3: the limit is a fixed point.} $T$ is continuous: given
	$\varepsilon>0$, take $\delta=\varepsilon$ (if $\beta=0$, any $\delta$); then
	$d(x,y)<\delta$ implies $d(Tx,Ty)\leq\beta d(x,y)<\varepsilon$. Hence, by
	\autoref{prop:sequential-continuity},
	\[
		Tx^\ast=T\Bigl(\lim_{n\to\infty}x_n\Bigr)=\lim_{n\to\infty}Tx_n
		=\lim_{n\to\infty}x_{n+1}=x^\ast .
	\]

	\emph{Step 4: uniqueness.} If $Tx^\ast=x^\ast$ and $Ty^\ast=y^\ast$, then
	$d(x^\ast,y^\ast)=d(Tx^\ast,Ty^\ast)\leq\beta\,d(x^\ast,y^\ast)$. If
	$d(x^\ast,y^\ast)>0$ this yields $1\leq\beta$, contradicting $\beta<1$. Hence
	$d(x^\ast,y^\ast)=0$ and $x^\ast=y^\ast$.

	\emph{Step 5: the error bound.} Let $m\to\infty$ in the Step 1 estimate
	$d(x_n,x_m)\leq\beta^n(1-\beta)^{-1}d(x_0,x_1)$. The metric is continuous in
	each argument --- an immediate consequence of the triangle inequality --- so the
	left-hand side tends to $d(x_n,x^\ast)$, giving
	\eqref{eq:contraction-error} with $x_1=Tx_0$.
\end{proof}

This is Theorem~8.3 of \textcite[p.~284]{fukuda2026notes}, proved over
\autocite[slides 64--72]{fukuda2026analysis}; the proof above follows those
steps, with the geometric sum and the passage to the limit in Step 5 written out.

\begin{remark}[What each hypothesis buys]\label{rem:contraction-hypotheses}
	Completeness is used once, in Step 2, and only to produce the limit; it is
	responsible for \emph{existence}. The strict inequality $\beta<1$ is used twice:
	in Step 1 to sum the geometric series, and in Step 4, where it is responsible for
	\emph{uniqueness}. The self-map property $T(X)\subseteq X$ is used implicitly
	throughout, in defining the iterates. Dropping any one of the three destroys a
	different part of the conclusion: on the complete space $\R$ the isometry
	$Tx=x+1$ ($\beta=1$) has no fixed point; on the incomplete space $(0,1]$ the
	contraction $Tx=x/2$ has none either, though its iterates are Cauchy.
\end{remark}

\begin{remark}[The error bound is what makes the theorem computational]
	\label{rem:error-bound}
	Inequality \eqref{eq:contraction-error} is an \emph{a priori} bound: it is
	computable from the first iteration alone, before the iteration is run. It
	converts the theorem from an existence statement into an algorithm with a stopping
	rule, since $n\geq\log\bigl(\varepsilon(1-\beta)/d(x_0,Tx_0)\bigr)/\log\beta$
	guarantees accuracy $\varepsilon$. Convergence is geometric at rate $\beta$, so
	the discount factor of a dynamic program is simultaneously an economic parameter
	and the convergence rate of the algorithm that solves it: a model with
	$\beta=0.99$ needs roughly $\log(0.5)/\log(0.99)\approx69$ times as many
	iterations per digit as one with $\beta=0.5$.
\end{remark}

\subsection{Three refinements}

\begin{corollary}[Iterates of a contraction]\label{cor:power-contraction}
	Let $X$ be a non-empty complete metric space and $T\colon X\to X$ a map such
	that $T^N$ is a contraction for some $N\in\N$. Then $T$ has a unique fixed
	point.
\end{corollary}

\begin{proof}
	By \autoref{thm:banach-fixed-point}, $T^N$ has a unique fixed point $x^\ast$.
	Then $T^N(Tx^\ast)=T(T^Nx^\ast)=Tx^\ast$, so $Tx^\ast$ is also a fixed point of
	$T^N$; by uniqueness $Tx^\ast=x^\ast$. Conversely any fixed point of $T$ is a
	fixed point of $T^N$, hence equals $x^\ast$.
\end{proof}

This is Exercise~8.30 of \textcite[p.~290]{fukuda2026notes}. It matters because
operators that are not contractions can have contracting powers: in a dynamic
program with periodic structure, or in a Markov chain whose transition matrix has
zero entries, the one-step operator may fail the modulus condition while the
$N$-step operator satisfies it.

\begin{corollary}[Properties inherited by the fixed point]\label{cor:invariant-set}
	Let $X$ be a complete metric space, $T\colon X\to X$ a contraction with fixed
	point $x^\ast$, and let $Y\subseteq X$ be non-empty and closed with
	$T(Y)\subseteq Y$. Then $x^\ast\in T(Y)\subseteq Y$.
\end{corollary}

\begin{proof}
	$Y$ is closed in a complete space, hence complete by
	\autoref{prop:closed-subset-complete}, and $T$ restricted to $Y$ is a
	contraction of $Y$ into itself. By \autoref{thm:banach-fixed-point} applied to
	$\restr{T}{Y}$ there is a fixed point in $Y$; by uniqueness in $X$ it equals
	$x^\ast$. Finally $x^\ast=Tx^\ast\in T(Y)$.
\end{proof}

\begin{remark}[How to prove that a value function is concave]\label{rem:property-P}
	\autoref{cor:invariant-set} is the standard route to qualitative properties of
	an object defined only as a fixed point. Let $P$ be a property, and let
	$Y=\{x\in X : x\text{ has }P\}$. If $Y$ is closed and $T$ preserves $P$, in the
	sense that $T(Y)\subseteq Y$, then the fixed point has $P$. In
	\autoref{ch:dynamic} this is applied three times over, with $X=\Cspace{S}{\R}$
	and $P$ successively ``increasing'', ``concave'' and ``strictly concave''. The
	closedness requirement is why the supremum norm is the right one:
	the pointwise limit of increasing functions is increasing and of concave
	functions is concave, and uniform convergence implies pointwise convergence, so
	both sets are closed in $\Bspace{S}{\R}$. Strict concavity is \emph{not} a
	closed property --- $f_k(s)=-s^2/k$ is strictly concave with uniform limit $0$
	--- and must be obtained by a separate argument
	\autocite[Remark 8.11, p.~290]{fukuda2026notes}.
\end{remark}

\begin{proposition}[Comparison of fixed points]\label{prop:fixed-point-comparison}
	Let $X$ be a complete metric space and let $S,T\colon X\to X$ be contractions
	with a common modulus $\beta\in[0,1)$ and fixed points $x_S,x_T$. If
	$d(Sx,Tx)\leq a$ for all $x\in X$, then
	\[
		d(x_S,x_T)\leq\frac{a}{1-\beta}.
	\]
\end{proposition}

\begin{proof}
	$d(x_S,x_T)=d(Sx_S,Tx_T)\leq d(Sx_S,Tx_S)+d(Tx_S,Tx_T)
		\leq a+\beta\,d(x_S,x_T)$. Rearranging and dividing by $1-\beta>0$ gives the
	bound.
\end{proof}

This is Exercise~8.29 of \textcite[p.~290]{fukuda2026notes}. It is the
stability result behind numerical dynamic programming: if the true Bellman
operator and its discretised approximation differ by at most $a$ uniformly, the
computed value function differs from the true one by at most $a/(1-\beta)$.
The factor $(1-\beta)^{-1}$ is an amplification, and it is large exactly when
the discount factor is close to $1$.

\section{Blackwell's Sufficient Conditions}\label{sec:comp-blackwell}

Verifying the contraction inequality directly for a Bellman operator requires
comparing two suprema, which is awkward. Blackwell's conditions replace that
comparison with two structural properties that can be read off the problem.

For $f,g\in\Bspace{S}{\R}$ write $f\leq g$ to mean $f(s)\leq g(s)$ for every
$s\in S$, and for $\alpha\in\R$ let $f+\alpha$ denote the function
$s\mapsto f(s)+\alpha$.

\begin{theorem}[Blackwell]\label{thm:blackwell}
	Let $S$ be a non-empty set and let $T\colon\Bspace{S}{\R}\to\Bspace{S}{\R}$
	satisfy
	\begin{enumerate}[(i)]
		\item \dfnas{monotonicity (Blackwell)}{monotonicity}: $f\leq g$ implies
		      $Tf\leq Tg$; and
		\item \dfnas{discounting (Blackwell)}{discounting}: there is
		      $\beta\in[0,1)$ such that
		      $T(f+\alpha)(s)\leq (Tf)(s)+\beta\alpha$ for every
		      $f\in\Bspace{S}{\R}$, every $\alpha\geq0$ and every $s\in S$.
	\end{enumerate}
	Then $T$ is a contraction with modulus $\beta$ on
	$\bigl(\Bspace{S}{\R},\norm{\cdot}_\infty\bigr)$, and therefore has a unique
	fixed point, reached from any starting point by iteration.
\end{theorem}

\begin{proof}
	Let $f,g\in\Bspace{S}{\R}$. For every $s$,
	$g(s)-f(s)\leq\abs{g(s)-f(s)}\leq\norm{f-g}_\infty$, so
	\[
		g\leq f+\norm{f-g}_\infty .
	\]
	Applying monotonicity and then discounting with
	$\alpha=\norm{f-g}_\infty\geq0$,
	\[
		Tg\;\leq\;T\bigl(f+\norm{f-g}_\infty\bigr)\;\leq\;Tf+\beta\norm{f-g}_\infty .
	\]
	Interchanging the roles of $f$ and $g$ gives
	$Tf\leq Tg+\beta\norm{f-g}_\infty$. The two inequalities together say that
	$\abs{(Tf)(s)-(Tg)(s)}\leq\beta\norm{f-g}_\infty$ for every $s\in S$, and
	taking the supremum over $s$ on the left yields
	$\norm{Tf-Tg}_\infty\leq\beta\norm{f-g}_\infty$. Since $\Bspace{S}{\R}$ is
	complete by \autoref{thm:B-complete}, \autoref{thm:banach-fixed-point} applies.
\end{proof}

This is Proposition~8.20 of the slides and Proposition~8.21 of
\textcite[p.~289]{fukuda2026notes}; the proof is the one sketched on
\autocite[slide 76]{fukuda2026analysis}.

\begin{remark}[The conditions are sufficient, not necessary]\label{rem:blackwell-sufficient}
	A contraction on $\Bspace{S}{\R}$ need not be monotone: on
	$\Bspace{S}{\R}$ with $S$ a single point, $Tf=-\tfrac12f$ is a contraction with
	modulus $\tfrac12$ and reverses order. What Blackwell's conditions provide is a
	criterion adapted to operators that arise from maximisation, for which both
	properties are transparent. Monotonicity of a Bellman operator says that a
	pointwise higher continuation value cannot lower today's optimised value, which
	holds because the same feasible choices remain available. Discounting says that
	adding a constant $\alpha$ to every continuation value raises today's value by
	at most $\beta\alpha$, which holds because the constant enters only through the
	discounted continuation term. Neither statement requires knowing the maximiser.
\end{remark}

\begin{eg}[Blackwell in $\R^n$]\label{eg:blackwell-rn}
	Equip $\R^n$ with $\norm{x}_\infty=\max_i\abs{x_i}$ and let
	$T\colon\R^n\to\R^n$ satisfy monotonicity ($x\leq y$ componentwise implies
	$Tx\leq Ty$) and discounting ($T(x+\alpha\mathbf1)\leq Tx+\beta\alpha\mathbf1$
	for $\alpha\geq0$, where $\mathbf1=(1,\dots,1)^\transp$). The proof of
	\autoref{thm:blackwell} applies verbatim with $S=\{1,\dots,n\}$, since
	$\Bspace{\{1,\dots,n\}}{\R}=\R^n$ with the supremum norm. This is
	Exercise~8.26 of \textcite[p.~290]{fukuda2026notes}, and it is the finite-state
	case: a Markov decision problem with $n$ states, where the Bellman operator is a
	map of $\R^n$ into itself and value iteration is a matrix--vector recursion.
\end{eg}

\section{Economic Applications}\label{sec:comp-applications}

\subsection{The Bellman operator}\label{subsec:comp-bellman}

\paragraph{Environment.} A planner observes a state $s\in S$, chooses an action
$a$ from a feasible set $\Gamma(s)\subseteq A$, receives a period return
$r(s,a)$, and moves to the state $g(s,a)$ next period. Future returns are
discounted at $\beta\in(0,1)$. The problem is
\begin{equation}\label{eq:sequence-problem}
	\sup_{(a_t)_{t\geq0}}\ \sum_{t=0}^{\infty}\beta^tr(s_t,a_t)
	\qquad\st\quad a_t\in\Gamma(s_t),\quad s_{t+1}=g(s_t,a_t),\quad s_0\text{ given}.
\end{equation}

\paragraph{The operator.} Define $T$ on $\Bspace{S}{\R}$ by
\begin{equation}\label{eq:bellman-operator}
	(Tv)(s)\eqdef\sup_{a\in\Gamma(s)}\Bigl\{r(s,a)+\beta\,v\bigl(g(s,a)\bigr)\Bigr\}.
\end{equation}

\begin{proposition}[The Bellman operator is a contraction]\label{prop:bellman-contraction}
	Suppose $r$ is bounded, $\Gamma(s)\neq\emptyset$ for every $s$, and
	$\beta\in[0,1)$. Then $T$ maps $\Bspace{S}{\R}$ into itself and is a contraction
	with modulus $\beta$; consequently \eqref{eq:bellman-operator} has a unique
	bounded solution $v^\ast$, and $T^nv_0\to v^\ast$ uniformly from any bounded
	$v_0$.
\end{proposition}

\begin{proof}
	\emph{Self-map.} If $\abs{r}\leq M_r$ and $\norm{v}_\infty\leq M_v$, then the
	expression in braces lies in $[-M_r-\beta M_v,\,M_r+\beta M_v]$ for every
	admissible $a$, so its supremum over the non-empty set $\Gamma(s)$ is finite and
	bounded by $M_r+\beta M_v$ uniformly in $s$. Hence $Tv\in\Bspace{S}{\R}$.

	\emph{Monotonicity.} If $v\leq w$ then for every $s$ and every
	$a\in\Gamma(s)$, $r(s,a)+\beta v(g(s,a))\leq r(s,a)+\beta w(g(s,a))$; taking
	suprema over the same set $\Gamma(s)$ preserves the inequality, so
	$Tv\leq Tw$.

	\emph{Discounting.} For $\alpha\geq0$,
	\[
		\bigl(T(v+\alpha)\bigr)(s)
		=\sup_{a\in\Gamma(s)}\bigl\{r(s,a)+\beta v(g(s,a))+\beta\alpha\bigr\}
		=(Tv)(s)+\beta\alpha,
	\]
	since an additive constant passes through a supremum. The required inequality
	holds with equality.

	Apply \autoref{thm:blackwell}.
\end{proof}

\paragraph{Scope of the fixed-point argument.} The theorem establishes that the
functional equation has exactly one bounded solution and that value iteration
converges to it geometrically, with the error after $n$ steps bounded by
\eqref{eq:contraction-error}. It does not by itself establish that $v^\ast$ is
the value of the sequence problem \eqref{eq:sequence-problem} --- that
equivalence requires a separate argument, given in \autoref{ch:dynamic} --- nor
that an optimal policy exists, which requires the supremum in
\eqref{eq:bellman-operator} to be attained and hence the Weierstrass theorem
(\autoref{thm:weierstrass}) applied to $\Gamma(s)$. Boundedness of $r$ is
restrictive: it excludes logarithmic utility on an unbounded state space, for
which a weighted supremum norm is used instead
\autocite[Sec.~4.3]{stokey1989recursive}.

\subsection{Asset pricing}\label{subsec:comp-lucas}

In the Lucas tree economy the price of a claim to a dividend stream satisfies
\begin{equation}\label{eq:lucas-operator}
	p(s)=\beta\,\E\bigl[\,m(s,s')\{d(s')+p(s')\}\mid s\,\bigr],
\end{equation}
where $s$ is the current state, $d$ the dividend and $m$ a positive stochastic
discount factor. Define $(Tp)(s)$ as the right-hand side. If
$\beta\,\sup_s\E[m(s,s')\mid s]\eqdef\bar\beta<1$ and $d$ is bounded, then for
bounded $p,q$,
\[
	\abs{(Tp)(s)-(Tq)(s)}
	=\beta\,\bigl|\E[m(s,s')\{p(s')-q(s')\}\mid s]\bigr|
	\leq\bar\beta\,\norm{p-q}_\infty,
\]
using positivity of $m$ and the conditional Jensen inequality of
\autoref{prop:conditional-jensen}. So $T$ is a contraction on
$\Bspace{S}{\R}$ and \eqref{eq:lucas-operator} has a unique bounded price
function. Uniqueness here is an economic statement: it rules out rational
bubbles among bounded price functions. Bubbles reappear precisely when the
boundedness restriction is dropped, since the homogeneous equation
$b(s)=\beta\E[m b(s')\mid s]$ then admits non-zero solutions
\autocite{lucas1978asset}.

\subsection{The steady-state Kalman filter}\label{subsec:comp-kalman}

\paragraph{Environment.} A hidden state follows
$\theta_{t+1}=\rho\,\theta_t+\eta_{t+1}$ with $\rho\in(-1,1)$ and
$\eta_{t+1}\sim N(0,\tau_2)$, and is observed each period through a signal
$y_t=\theta_t+\varepsilon_t$ with $\varepsilon_t\sim N(0,\tau_1^{-1})$, all
shocks independent. Following the source exercise, $\tau_2$ is the
\emph{variance} of the state innovation and $\tau_1$ the \emph{precision} of the
signal. An observer forms Bayesian beliefs about $\theta_t$ from the signal
history; let $x_t$ denote the variance of that belief \emph{before} seeing
$y_t$.

\paragraph{Recursion.} Normal--normal updating (\autoref{thm:normal-updating})
adds precisions: after observing $y_t$ the posterior variance is
$\bigl(x_t^{-1}+\tau_1\bigr)^{-1}$. The transition then scales the variance by
$\rho^2$ and adds the innovation variance $\tau_2$, so $x_{t+1}=T(x_t)$ with
\begin{equation}\label{eq:riccati}
	T(x)=\rho^2\bigl(x^{-1}+\tau_1\bigr)^{-1}+\tau_2
	=\frac{\rho^2x}{1+\tau_1x}+\tau_2,
	\qquad x\in\R_{++} .
\end{equation}
This is the scalar Riccati equation of the Kalman filter.

\begin{proposition}[The Riccati operator is a contraction]\label{prop:riccati-contraction}
	Let $\rho\in(-1,1)$ and $\tau_1,\tau_2>0$. Then $T$ in \eqref{eq:riccati} maps
	$\R_{++}$ into itself and is a contraction with modulus $\rho^2\in[0,1)$.
	Consequently the conditional variance converges, from any initial belief, to the
	unique fixed point
	\begin{equation}\label{eq:riccati-fixed-point}
		x^\ast=\frac{(\rho^2+\tau_1\tau_2-1)+\sqrt{(1-\rho^2-\tau_1\tau_2)^2+4\tau_1\tau_2}}{2\tau_1}\,,
	\end{equation}
	the unique positive root of $\tau_1x^2+(1-\rho^2-\tau_1\tau_2)x-\tau_2=0$.
\end{proposition}

\begin{proof}
	\emph{Self-map.} For $x>0$ both terms of $\rho^2x/(1+\tau_1x)+\tau_2$ are
	non-negative and the second is strictly positive, so $T(x)\geq\tau_2>0$.

	\emph{Contraction.} For $x,y>0$,
	\[
		T(x)-T(y)=\rho^2\left(\frac{x}{1+\tau_1x}-\frac{y}{1+\tau_1y}\right)
		=\frac{\rho^2(x-y)}{(1+\tau_1x)(1+\tau_1y)},
	\]
	after clearing the common denominator and cancelling the terms in $\tau_1xy$.
	Since $\tau_1>0$ and $x,y>0$, the denominator exceeds $1$, so
	$\abs{T(x)-T(y)}\leq\rho^2\abs{x-y}$ with $\rho^2<1$. The set
	$[\tau_2,\rho^2\tau_1^{-1}+\tau_2]$ is closed, non-empty, and invariant under
	$T$, hence complete by \autoref{prop:closed-subset-complete};
	\autoref{thm:banach-fixed-point} applies on it, and
	\autoref{cor:invariant-set} places the fixed point there.

	\emph{Fixed point.} Setting $x=\rho^2x/(1+\tau_1x)+\tau_2$ and multiplying by
	$1+\tau_1x>0$ gives $x+\tau_1x^2=\rho^2x+\tau_2+\tau_1\tau_2x$, that is
	$\tau_1x^2+(1-\rho^2-\tau_1\tau_2)x-\tau_2=0$. The product of the two roots is
	$-\tau_2/\tau_1<0$, so exactly one root is positive, and it is
	\eqref{eq:riccati-fixed-point}. As a check, $\rho=0$ gives
	$\tau_1x^2+(1-\tau_1\tau_2)x-\tau_2=(\tau_1x+1)(x-\tau_2)=0$, hence
	$x^\ast=\tau_2$: with no persistence, yesterday's signal is worthless and the
	conditional variance is the innovation variance.
\end{proof}

\paragraph{Economic content.} First, uncertainty neither vanishes
nor explodes. However long the observer has been learning, the conditional
variance settles at $x^\ast>0$: new innovations arrive at exactly the rate at
which signals resolve old ones. Permanent learning does not produce perfect
information, which is why models of learning about a persistent fundamental
sustain non-degenerate belief dispersion indefinitely.

Second, the rate of convergence is $\rho^2$, the square of the persistence.
Beliefs about a near-random-walk state ($\rho\to1$) converge to their steady
state arbitrarily slowly, so the transitional dynamics of learning are
economically relevant over long horizons precisely when the fundamental is
persistent.

Third, comparative statics are read off
\eqref{eq:riccati-fixed-point} through the quadratic. Write
$F(x)\eqdef\tau_1x^2+(1-\rho^2-\tau_1\tau_2)x-\tau_2$, so that $F(x^\ast)=0$ and
$F'(x^\ast)>0$ at the larger root. Since
$\partial F/\partial\rho^2=-x^\ast<0$, the implicit function theorem
(\autoref{thm:implicit-function}) gives
$\partial x^\ast/\partial\rho^2=x^\ast/F'(x^\ast)>0$: a more persistent state is
harder to track, not easier. Likewise
$\partial F/\partial\tau_1=x^\ast(x^\ast-\tau_2)$, which is strictly positive
because $x^\ast=\rho^2x^\ast/(1+\tau_1x^\ast)+\tau_2>\tau_2$; hence
$\partial x^\ast/\partial\tau_1<0$ and a more precise signal lowers
steady-state uncertainty, without qualification. This is Exercise~8.27 of
\textcite[p.~290]{fukuda2026notes}, with the interpretation supplied.

\subsection{Input--output systems revisited}\label{subsec:comp-leontief}

\autoref{eg:leontief} solved $x=Ax+d$ by a Neumann series under
$\rho(A)<1$. \autoref{thm:banach-fixed-point} gives the same conclusion with a
different argument and a computable error bound: if $\norm{A}<1$ in some
operator norm, then $Tx=Ax+d$ satisfies
$\norm{Tx-Ty}=\norm{A(x-y)}\leq\norm{A}\norm{x-y}$, so $T$ is a contraction on
the complete space $\R^n$ and the iteration $x_{k+1}=Ax_k+d$ converges to the
unique solution from any starting point, with
$\norm{x_k-x^\ast}\leq\norm{A}^k(1-\norm{A})^{-1}\norm{x_1-x_0}$. The two
hypotheses are not identical: $\rho(A)\leq\norm{A}$ always, so the spectral
condition is weaker, and by \autoref{cor:power-contraction} the gap is exactly
accounted for by passing to a power of $T$, since $\rho(A)<1$ implies
$\norm{A^N}<1$ for some $N$.

\section{Chapter Summary}\label{sec:comp-summary}

A Cauchy sequence is one whose terms cluster without reference to any limit; every
convergent sequence is Cauchy and every Cauchy sequence is bounded, but the
converse of the first implication holds only in complete spaces. $\R^n$ is
complete, as a consequence of the least upper bound axiom via
Bolzano--Weierstrass; a subset of a complete space is complete exactly when it is
closed (\autoref{prop:closed-subset-complete}); the bounded functions on any set
form a Banach space under the supremum norm (\autoref{thm:B-complete}), and the
bounded continuous functions form a closed subspace of it.

A contraction on a non-empty complete metric space has exactly one fixed point,
reached by iteration from any starting point at geometric rate $\beta$ with the
computable error bound \eqref{eq:contraction-error}. Completeness delivers
existence and $\beta<1$ delivers uniqueness; each is indispensable. Three
refinements extend the reach of the theorem: a map with a contracting power has a
unique fixed point; a closed invariant set contains the fixed point, which is how
qualitative properties of value functions are established; and fixed points of
uniformly close contractions are within $a/(1-\beta)$ of each other, which bounds
the error of numerical approximation. Blackwell's monotonicity and discounting
conditions are sufficient for the contraction property and are verifiable
directly on operators defined by maximisation, which is what makes them the
standard tool for Bellman operators, asset-pricing operators and Bayesian
updating recursions.

\begin{exercise}\label{ex:comp-1}
	Let $X$ be a complete normed space and $T\colon X\to X$ a contraction with
	modulus $\beta$. Show that for every $y\in X$ there is a unique $x\in X$ with
	$Tx-x=y$, and that the solution map $y\mapsto x$ is Lipschitz with constant
	$(1-\beta)^{-1}$. Interpret the second statement as a stability property of the
	linear system $x=Ax+d$ with respect to perturbations of $d$
	\autocite[Ex.~8.31]{fukuda2026notes}.
\end{exercise}

\begin{exercise}\label{ex:comp-2}
	Show that $Tx=x+1/x$ on $[1,\infty)$ satisfies $\abs{Tx-Ty}<\abs{x-y}$ for all
	$x\neq y$ and has no fixed point, and that $Tx=\sqrt{1+x^2}$ on $\R$ does the
	same. Identify the supremum of $\abs{Tx-Ty}/\abs{x-y}$ in each case and explain
	why \autoref{thm:banach-fixed-point} does not apply.
\end{exercise}

\begin{exercise}\label{ex:comp-3}
	Let $S=\{1,\dots,n\}$ and let $Q$ be an $n\times n$ row-stochastic matrix. Show
	that $Tv=r+\beta Qv$ is a contraction on $(\R^n,\norm{\cdot}_\infty)$ with
	modulus $\beta$, that its fixed point is $v^\ast=(I-\beta Q)^{-1}r$, and that
	$I-\beta Q$ is invertible for every $\beta\in[0,1)$. Verify Blackwell's two
	conditions for $T$ directly, and explain which property of $Q$ delivers each.
\end{exercise}
